% R8_Domain_Folding.tex
% monogate Cost Theory — R8: Domain-Folding Theorem
% Proves exact savings from positive-domain routing in SuperBEST v3,
% with static analysis rules and catalog-wide quantification.
% Date: 2026-04-20
% Author: Arturo R. Almaguer
% Classification: THEOREM

\documentclass[12pt,a4paper]{article}
\usepackage{amsmath,amssymb,amsthm}
\usepackage{geometry}
\usepackage{booktabs}
\usepackage{array}
\usepackage{longtable}
\usepackage{xcolor}
\usepackage{hyperref}
\usepackage{listings}
\geometry{margin=2.5cm}

% ── Theorem environments ─────────────────────────────────────────────────────
\newtheorem{theorem}{Theorem}[section]
\newtheorem{lemma}[theorem]{Lemma}
\newtheorem{corollary}[theorem]{Corollary}
\newtheorem{proposition}[theorem]{Proposition}
\theoremstyle{definition}
\newtheorem{definition}[theorem]{Definition}
\newtheorem{conjecture}[theorem]{Conjecture}
\newtheorem{remark}[theorem]{Remark}
\newtheorem{example}[theorem]{Example}
\newtheorem{observation}[theorem]{Observation}
\newtheorem{algorithm}[theorem]{Algorithm}

% ── Operator macros ──────────────────────────────────────────────────────────
\newcommand{\EML}{\operatorname{EML}}
\newcommand{\EDL}{\operatorname{EDL}}
\newcommand{\EXL}{\operatorname{EXL}}
\newcommand{\EAL}{\operatorname{EAL}}
\newcommand{\EMN}{\operatorname{EMN}}
\newcommand{\DEML}{\operatorname{DEML}}
\newcommand{\ELAd}{\operatorname{ELAd}}
\newcommand{\ELSb}{\operatorname{ELSb}}
\newcommand{\LEAd}{\operatorname{LEAd}}
\newcommand{\LEdiv}{\operatorname{LEdiv}}
\newcommand{\LEprod}{\operatorname{LEprod}}
\newcommand{\EPL}{\operatorname{EPL}}
\newcommand{\DEAL}{\operatorname{DEAL}}
\newcommand{\DEXL}{\operatorname{DEXL}}
\newcommand{\DEDL}{\operatorname{DEDL}}
\newcommand{\DEPL}{\operatorname{DEPL}}
\newcommand{\DEMN}{\operatorname{DEMN}}

% ── Cost macros ───────────────────────────────────────────────────────────────
\newcommand{\Cost}{\operatorname{Cost}}
\newcommand{\NC}{\operatorname{NaiveCost}}
\newcommand{\SD}{\operatorname{SharingDiscount}}
\newcommand{\PB}{\operatorname{PatternBonus}}
\newcommand{\Fam}{\mathcal{F}}

% ── Domain macros ─────────────────────────────────────────────────────────────
\newcommand{\POS}{\textsc{pos}}
\newcommand{\NEG}{\textsc{neg}}
\newcommand{\UNK}{\textsc{unk}}

% ── Column types ──────────────────────────────────────────────────────────────
\newcolumntype{L}[1]{>{\raggedright\arraybackslash}p{#1}}
\newcolumntype{C}[1]{>{\centering\arraybackslash}p{#1}}

% ── Listings style ────────────────────────────────────────────────────────────
\lstset{
  basicstyle=\ttfamily\small,
  keywordstyle=\bfseries,
  commentstyle=\itshape\color{gray},
  frame=single,
  xleftmargin=1em,
  xrightmargin=1em,
  breaklines=true,
  columns=flexible,
  keepspaces=true,
}

\title{The Domain-Folding Theorem\\[6pt]
       \large T\textsubscript{DF}: Exact Savings from Positive-Domain Routing
       in SuperBEST~v3,\\
       Static Analysis Rules, and Catalog-Wide Quantification}
\author{Arturo R.\ Almaguer\\
\small monogate research \quad \texttt{almaguer1986@gmail.com}}
\date{April 2026}

% ═════════════════════════════════════════════════════════════════════════════
\begin{document}
\maketitle

% ── Abstract ─────────────────────────────────────────────────────────────────
\begin{abstract}
\textbf{Historical note (SuperBEST v5, ADD-T1, April 2026).}
Theorem ADD-T1 proves that general-domain addition costs exactly 2 nodes via
$\LEdiv(x,\DEML(y,1)) = x+y$, eliminating the positive/general routing split
entirely.  Both routes now cost 2 nodes; the domain-folding saving $\Delta_{\mathrm{add}} = 0$.
The analysis below is preserved for historical completeness and as an illustration
of domain propagation techniques applicable to other operator families.
\bigskip

SuperBEST v3 provides two routing paths for binary addition: a positive-domain
route costing $3n$ nodes, used when both operands are provably positive, and a
general-domain route costing $11n$ nodes, used otherwise (requiring internal
absolute-value circuitry).
The \emph{Domain-Folding Theorem} ($T_{\mathrm{DF}}$) formalises exactly when the
positive-domain route applies and quantifies the savings: \emph{8 nodes per
addition} (v3/v4; \emph{0 nodes per addition in v5}).
We prove eight domain propagation rules that a static analyser can apply to an
expression tree to classify operands as $\POS$, $\NEG$, or $\UNK$ without
evaluating any inputs.
Algorithm~2 (\textsc{PositiveDomainAnalysis}) implements this traversal in
linear time.
Applied to the 157-equation monogate catalog, we estimate that approximately
60\% of equations contain at least one positively-routable addition, yielding a
total savings of roughly $8n \times 151$ addition sites across the catalog ---
a non-trivial structural economy over the na\"{i}ve general-domain assumption
(relevant for v3/v4; superseded by ADD-T1 in v5).
\end{abstract}

\tableofcontents
\bigskip

% ═════════════════════════════════════════════════════════════════════════════
\section{Background: The 8-Node Saving Per Addition}
\label{sec:background}
% ═════════════════════════════════════════════════════════════════════════════

\subsection{SuperBEST v3 Addition Costs}

The monogate SuperBEST v3 routing table (reproduced in
Table~\ref{tab:add-costs}) distinguishes two implementation paths for binary
addition $x + y$.

\begin{table}[h]
\centering
\caption{Addition costs by domain (v3/v4 historical; v5 unified).}
\label{tab:add-costs}
\smallskip
\begin{tabular}{@{}lccc@{}}
\toprule
\textbf{Routing path} & \textbf{Condition} & \textbf{Cost v3/v4} & \textbf{Cost v5} \\
\midrule
Positive-domain addition & $x > 0$ and $y > 0$ provably & $3n$ & $2n$ \\
General-domain addition  & at least one operand sign unknown & $11n$ & $2n$ \\
\bottomrule
\end{tabular}
\smallskip

\noindent\textbf{v5 note:} ADD-T1 proves \texttt{add}$(x,y) = 2n$ for all real $x,y$.
The routing distinction is eliminated.
\end{table}

\begin{definition}[Domain routing]
\label{def:domain-routing}
Let $E = \mathrm{add}(E_1, E_2)$ be a binary addition node in an EML expression
tree.  The \emph{domain routing decision} is the compile-time selection between:
\begin{itemize}
  \item \textbf{Positive route}: chosen when a static analysis certifies
        $E_1 > 0$ and $E_2 > 0$ for all valid inputs; costs $3n$ nodes.
  \item \textbf{General route}: chosen in all other cases; costs $11n$ nodes
        (the extra $8n$ implement an absolute-value sub-circuit that handles
        negative operands correctly).
\end{itemize}
\end{definition}

\begin{definition}[Saving per addition]
\label{def:saving}
The \emph{domain-folding saving} for a single addition node is
\[
  \Delta_{\mathrm{add}} \;=\; 11n - 3n \;=\; 8n \quad(\text{v3/v4}).
\]
This saving is realised if and only if both operands are certified positive by
the static analyser.  \emph{In SuperBEST v5 (ADD-T1), $\Delta_{\mathrm{add}} = 2n - 2n = 0$:
all additions cost $2n$ regardless of domain, so domain folding for addition
is obsolete.}
\end{definition}

\subsection{Motivation and Scope}

\textbf{Historical context (v3/v4).}
The distinction between $3n$ and $11n$ was the largest single-operator cost gap
in $\mathcal{F}_{16}$ (compare: multiplication costs $2n$ regardless of sign,
and \texttt{exp} costs $1n$ regardless of sign).
In SuperBEST v5, this gap is eliminated by ADD-T1: \texttt{add} = $2n$ uniformly.
Consequently, domain-aware compilation has a larger per-site impact on total
cost than any other routing decision available in SuperBEST v3.

The key question is: for how many addition sites in scientific equations can we
certify positivity at compile time, without knowing the actual input values?
The domain propagation rules in Section~\ref{sec:rules} and the static analysis
algorithm in Section~\ref{sec:algorithm} answer this question systematically.

% ═════════════════════════════════════════════════════════════════════════════
\section{Positive-Domain Propagation Rules}
\label{sec:rules}
% ═════════════════════════════════════════════════════════════════════════════

We prove eight rules that a static analyser can apply bottom-up through an
expression tree to derive provable domain labels ($\POS$, $\NEG$, or $\UNK$)
for every sub-expression.

\begin{definition}[Domain label]
\label{def:domain-label}
A \emph{domain label} $\delta(E)$ is assigned to an expression $E$ and takes
one of three values:
\begin{align*}
  \delta(E) = \POS &\quad\Longleftrightarrow\quad E(x) > 0 \text{ for all valid inputs } x, \\
  \delta(E) = \NEG &\quad\Longleftrightarrow\quad E(x) < 0 \text{ for all valid inputs } x, \\
  \delta(E) = \UNK &\quad\Longleftrightarrow\quad \text{sign of } E \text{ cannot be certified statically.}
\end{align*}
\end{definition}

\begin{lemma}[Rule R1: Exponential is always positive]
\label{lem:R1}
For any expression $A$ over $\mathbb{R}$,
\[
  \delta(\exp(A)) \;=\; \POS.
\]
\end{lemma}

\begin{proof}
By definition, $\exp(A) = e^{A}$.  The exponential function $e^t$ satisfies
$e^t > 0$ for all $t \in \mathbb{R}$ (since $e^t = \lim_{n\to\infty}(1+t/n)^n$
and the sequence is bounded below by $e^{-|t|} > 0$; equivalently, $e^t \cdot
e^{-t} = e^0 = 1 > 0$ with $e^{-t} > 0$ implying $e^t > 0$).
Hence $\delta(\exp(A)) = \POS$ regardless of $\delta(A)$.
\end{proof}

\begin{lemma}[Rule R2: Square is non-negative]
\label{lem:R2}
For any expression $A$ over $\mathbb{R}$,
\[
  A^2 \;\geq\; 0,
\]
with equality if and only if $A = 0$.  In particular, for any $A$ with
$\delta(A) \neq \{0\}$ (i.e., $A$ is not identically zero), $\delta(A^2) =
\POS$.
\end{lemma}

\begin{proof}
For any real number $t$, $t^2 \geq 0$ by the ordered-field property of
$\mathbb{R}$: the product of a number with itself has the same sign as $t \cdot t
= |t|^2 \geq 0$.  Equality holds iff $t = 0$.  When the static analyser knows
that $A$ is not identically zero (e.g., $A$ is a free variable over a positive
domain, or $A = \exp(B)$ which is non-zero), $A^2 > 0$ certifiably.
\end{proof}

\begin{lemma}[Rule R3: Absolute value is non-negative]
\label{lem:R3}
For any expression $A$ over $\mathbb{R}$,
\[
  |A| \;\geq\; 0.
\]
When $A$ is not identically zero, $|A| > 0$ and $\delta(|A|) = \POS$.
\end{lemma}

\begin{proof}
By definition, $|A| = \max(A, -A) \geq 0$.  Equality holds iff $A = 0$.
For non-zero $A$, the absolute value is strictly positive.
\end{proof}

\begin{lemma}[Rule R4: Product of positives is positive]
\label{lem:R4}
If $\delta(E_1) = \POS$ and $\delta(E_2) = \POS$, then
\[
  \delta(E_1 \cdot E_2) \;=\; \POS.
\]
If $\delta(E_1) = \NEG$ and $\delta(E_2) = \NEG$, then
$\delta(E_1 \cdot E_2) = \POS$.
\end{lemma}

\begin{proof}
In an ordered field, the product of two positive numbers is positive: if $a > 0$
and $b > 0$, then $ab > 0$.  Similarly, if $a < 0$ and $b < 0$, then
$(-a)(-b) = ab > 0$.  Both statements follow directly from the ordered-field
axioms.
\end{proof}

\begin{lemma}[Rule R5: Quotient of positives is positive]
\label{lem:R5}
If $\delta(E_1) = \POS$ and $\delta(E_2) = \POS$, then
\[
  \delta(E_1 / E_2) \;=\; \POS.
\]
\end{lemma}

\begin{proof}
If $a > 0$ and $b > 0$, then $a/b = a \cdot (1/b)$.  Since $b > 0$ we have
$1/b > 0$, and by Rule~R4, $a \cdot (1/b) > 0$.
\end{proof}

\begin{lemma}[Rule R6: Logarithm of a positive is real]
\label{lem:R6}
If $\delta(E) = \POS$, then $\ln(E)$ is well-defined (real-valued), but
$\delta(\ln(E)) = \UNK$ in general.
\end{lemma}

\begin{proof}
The natural logarithm $\ln : (0, \infty) \to \mathbb{R}$ is defined for all
positive arguments, so $\delta(E) = \POS$ guarantees domain validity.
However, $\ln(E)$ ranges over all of $\mathbb{R}$: $\ln(E) > 0$ when $E > 1$,
$\ln(E) = 0$ when $E = 1$, and $\ln(E) < 0$ when $0 < E < 1$.  Without further
constraints on the magnitude of $E$, the sign of $\ln(E)$ cannot be certified
statically.  Hence $\delta(\ln(E)) = \UNK$.
\end{proof}

\begin{remark}
Rule R6 shows that $\ln$ is a domain-sign ``sink'': even when the input is
provably positive, the output sign is unknown.  This is in sharp contrast to
$\exp$, which is a domain-sign ``source'' (always produces $\POS$).
\end{remark}

\begin{lemma}[Rule R7: Positive scalar multiple]
\label{lem:R7}
If $\delta(E) = \POS$ and $c > 0$ is a positive compile-time constant, then
\[
  \delta(c \cdot E) \;=\; \POS.
\]
\end{lemma}

\begin{proof}
This is a special case of Rule~R4 with $E_1 = c$ (a positive constant, so
$\delta(c) = \POS$) and $E_2 = E$.
\end{proof}

\begin{lemma}[Rule R8: Positive base to positive exponent]
\label{lem:R8}
If $\delta(E) = \POS$ and $c > 0$ is a positive constant, then
\[
  \delta(E^c) \;=\; \POS.
\]
\end{lemma}

\begin{proof}
For $e > 0$ and $c > 0$, $e^c = \exp(c \ln e)$.  Since $e > 0$, $\ln e$ is
real-valued (Rule~R6), and $c \cdot \ln e$ is real, so $\exp(c \ln e) > 0$ by
Rule~R1.  Alternatively, $e^c > 0$ follows directly from the fact that any
positive real raised to a real power remains positive: the function $t \mapsto
e^c$ is continuous and equals $e^c = \lim_{n} r_n^c$ where $r_n \to e$,
$r_n > 0$, and each $r_n^c > 0$.
\end{proof}

\subsection{Domain Propagation Table}
\label{sec:prop-table}

Table~\ref{tab:propagation} summarises the output domain label for each operator
in $\mathcal{F}_{16}$ as a function of input domain labels.

\begin{longtable}{@{}L{3.2cm}L{3.2cm}C{2.5cm}L{4.5cm}@{}}
\caption{Domain propagation table for $\mathcal{F}_{16}$ operators.
         ``P'' = $\POS$, ``N'' = $\NEG$, ``?'' = $\UNK$.}
\label{tab:propagation}\\
\toprule
\textbf{Operator} & \textbf{Input domain(s)} & \textbf{Output domain} &
\textbf{Rule / justification} \\
\midrule
\endfirsthead
\multicolumn{4}{l}{\small\itshape (continued)} \\
\toprule
\textbf{Operator} & \textbf{Input domain(s)} & \textbf{Output domain} &
\textbf{Rule / justification} \\
\midrule
\endhead
\bottomrule
\endfoot
$\exp(A)$        & any                   & P   & R1 \\[3pt]
$A^2$            & any (non-zero)        & P   & R2 \\[3pt]
$|A|$            & any (non-zero)        & P   & R3 \\[3pt]
$A \cdot B$      & P, P                  & P   & R4 \\
                 & N, N                  & P   & R4 \\
                 & P, N or N, P          & N   & R4 \\
                 & ?, any or any, ?      & ?   & cannot certify \\[3pt]
$A / B$          & P, P                  & P   & R5 \\
                 & N, N                  & P   & R5 \\
                 & P, N or N, P          & N   & \\
                 & ?, any or any, ?      & ?   & \\[3pt]
$\ln(A)$         & P                     & ?   & R6 \\
                 & N or ?                & undefined / ? & domain error if N \\[3pt]
$c \cdot A$      & $c>0$, P              & P   & R7 \\
                 & $c>0$, N              & N   & R7 \\
                 & $c<0$, P              & N   & R7 \\
                 & $c<0$, N              & P   & R7 \\[3pt]
$A^c$ ($c>0$)    & P                     & P   & R8 \\
                 & ? or N (real $c$)     & ?   & may not be real \\[3pt]
$\mathrm{add}(A,B)$ & P, P              & P, \emph{use pos-route} & T\textsubscript{DF} \\
                 & any other             & ?   & use gen-route \\[3pt]
$\mathrm{neg}(A)$ & P                   & N   & sign flip \\
                  & N                   & P   & sign flip \\
                  & ?                   & ?   & \\[3pt]
$\mathrm{sub}(A,B)$ & any              & ?   & sign not certifiable \\[3pt]
constant $c$      & ---                 & P if $c>0$, N if $c<0$, ? if $c=0$ & base case \\[3pt]
variable $x$      & ---                 & from input declaration & base case \\
\end{longtable}

% ═════════════════════════════════════════════════════════════════════════════
\section{Static Analysis Algorithm}
\label{sec:algorithm}
% ═════════════════════════════════════════════════════════════════════════════

\begin{definition}[Domain declaration]
\label{def:domain-decl}
A \emph{domain declaration} is a compile-time annotation attaching a domain
label to a free variable: e.g., $x : \POS$ declares that $x > 0$ for all valid
inputs.  Domain declarations are part of the expression metadata; they are not
evaluated at run time.
\end{definition}

\begin{algorithm}[\textsc{PositiveDomainAnalysis}]
\label{alg:PDA}
Algorithm~2 traverses an expression tree $E$ bottom-up, assigning a domain
label $\delta(E)$ to each node and flagging each \texttt{add} node as
\emph{positive-routed} (PR) or \emph{general-routed} (GR).
\end{algorithm}

\begin{lstlisting}[caption={Algorithm 2: PositiveDomainAnalysis(E)},
                   label={lst:PDA}]
-- Input:  expression tree E,
--         domain declarations D : Var -> {POS, NEG, UNK}
-- Output: delta(E) in {POS, NEG, UNK};
--         routing label for each add node (PR or GR)

function domain_analyze(E, D):

  -- Base cases -------------------------------------------------------
  if E is a numeric constant c:
    if   c > 0: return POS
    elif c < 0: return NEG
    else:       return UNK   -- c = 0

  if E is a free variable x:
    return D(x)              -- from domain declarations

  -- Unary operators --------------------------------------------------
  if E = exp(A):
    domain_analyze(A, D)     -- recurse (side effects for add nodes)
    return POS               -- R1: always positive

  if E = neg(A):
    dA = domain_analyze(A, D)
    if   dA = POS: return NEG
    elif dA = NEG: return POS
    else:          return UNK

  if E = abs(A):
    domain_analyze(A, D)
    return POS               -- R3: always non-negative

  if E = ln(A):
    dA = domain_analyze(A, D)
    -- domain validity check (log requires positive input)
    if dA = NEG: raise DomainError("ln of negative")
    return UNK               -- R6: real but sign unknown

  if E = pow(A, c):          -- A^c, c a positive constant
    dA = domain_analyze(A, D)
    if dA = POS: return POS  -- R8
    return UNK

  -- Binary operators -------------------------------------------------
  if E = mul(A, B):
    dA = domain_analyze(A, D)
    dB = domain_analyze(B, D)
    if   dA = POS and dB = POS: return POS
    elif dA = NEG and dB = NEG: return POS
    elif dA = POS and dB = NEG: return NEG
    elif dA = NEG and dB = POS: return NEG
    else: return UNK

  if E = div(A, B):
    dA = domain_analyze(A, D)
    dB = domain_analyze(B, D)
    if   dA = POS and dB = POS: return POS
    elif dA = NEG and dB = NEG: return POS
    elif dA = POS and dB = NEG: return NEG
    elif dA = NEG and dB = POS: return NEG
    else: return UNK

  if E = add(A, B):
    dA = domain_analyze(A, D)
    dB = domain_analyze(B, D)
    if dA = POS and dB = POS:
      label(E, PR)           -- positive-domain route, cost 3n
      return POS
    else:
      label(E, GR)           -- general-domain route, cost 11n
      if dA = POS and dB = NEG: return UNK
      if dA = NEG and dB = POS: return UNK
      if dA = NEG and dB = NEG: return NEG
      return UNK

  if E = sub(A, B):
    dA = domain_analyze(A, D)
    dB = domain_analyze(B, D)
    -- subtraction: A - B = A + (-B); sign not certifiable in general
    return UNK
\end{lstlisting}

\begin{proposition}[Correctness of Algorithm~2]
\label{prop:correctness}
Let $E$ be an expression tree and $D$ a domain declaration for all free variables
of $E$.  Algorithm~\ref{alg:PDA} runs in $O(|E|)$ time (where $|E|$ is the number
of nodes in $E$) and satisfies the following soundness guarantee: if the
algorithm assigns $\delta(E) = \POS$, then $E(x) > 0$ for all inputs $x$
consistent with $D$.
\end{proposition}

\begin{proof}
\textbf{Time complexity.}
Each node of $E$ is visited exactly once in the bottom-up traversal.  The work
per node is $O(1)$ (one conditional branch on the operator type).  Hence total
time is $O(|E|)$.

\textbf{Soundness.}
We prove by structural induction on $E$ that whenever the algorithm returns
$\POS$, the expression is genuinely positive.
\begin{itemize}
  \item \emph{Base: constant $c > 0$.}  $c > 0$ by the branch condition.
  \item \emph{Base: variable $x$ with $D(x) = \POS$.}  By the semantics of
        domain declarations, $x > 0$ for all valid inputs.
  \item \emph{Inductive step.}  For each operator, the returned $\POS$ label
        follows from one of Lemmas~\ref{lem:R1}--\ref{lem:R8}, whose proofs
        rely only on the ordered-field axioms of $\mathbb{R}$.  The inductive
        hypothesis gives that the input labels are sound, and each lemma
        propagates soundness to the output.
\end{itemize}
The algorithm never returns $\POS$ without invoking one of the eight rules, so
soundness holds globally.

\textbf{No false negatives (incompleteness note).}
The algorithm may assign $\UNK$ to expressions that are in fact positive (e.g.,
$\ln(e^x) = x$ which may be positive, or $x - (-|x|)$ which simplifies to a
non-negative).  This incompleteness is intentional: the algorithm is
\emph{conservative}, never unsafely applying the cheaper route.
\end{proof}

\begin{corollary}[Routing assignment is safe]
\label{cor:routing-safe}
Every addition node labelled $\mathrm{PR}$ (positive route) by Algorithm~\ref{alg:PDA}
is correctly routed at cost $3n$.  No correctness error can arise from applying
the positive-domain circuit to an addition whose operands are actually
non-positive, because the algorithm only labels $\mathrm{PR}$ when both
operands are certified $\POS$.
\end{corollary}

% ═════════════════════════════════════════════════════════════════════════════
\section{The Domain-Folding Theorem}
\label{sec:theorem}
% ═════════════════════════════════════════════════════════════════════════════

\begin{theorem}[$T_{\mathrm{DF}}$ --- Domain-Folding Theorem]
\label{thm:TDF}
Let $E_1$ and $E_2$ be expressions in an EML tree such that
$\delta(E_1) = \POS$ and $\delta(E_2) = \POS$ (as certified by
Algorithm~\ref{alg:PDA} under the input domain declarations $D$).
Then:
\begin{enumerate}
  \item $\Cost(\mathrm{add}(E_1, E_2))$ uses the positive-domain routing at
        $3n$ nodes, not the general-domain routing at $11n$ nodes.
  \item The savings per such addition is exactly
        \[
          \Delta_{\mathrm{add}} \;=\; 11n - 3n \;=\; 8n.
        \]
  \item The total savings for an expression $E$ with $k$ positively-routable
        addition nodes (as identified by Algorithm~\ref{alg:PDA}) is exactly
        \[
          \Delta(E) \;=\; 8kn.
        \]
\end{enumerate}
\end{theorem}

\begin{proof}
\textbf{Part (1).}
By Definition~\ref{def:domain-routing}, the positive route is applied if and
only if both operands are certified positive.
Algorithm~\ref{alg:PDA} returns $\POS$ for $E_1$ and $E_2$ only when
Lemmas~\ref{lem:R1}--\ref{lem:R8} guarantee positivity for all valid inputs.
By Proposition~\ref{prop:correctness} (soundness), these guarantees are
correct.  Hence the positive-domain circuit is applicable, and the SuperBEST v3
router selects it at cost $3n$.

\textbf{Part (2).}
The general-domain route costs $11n$ (Table~\ref{tab:add-costs}).
The positive-domain route costs $3n$ (Table~\ref{tab:add-costs}).
The difference is $11 - 3 = 8$ nodes: $\Delta_{\mathrm{add}} = 8n$.

\textbf{Part (3).}
Let $a_1, \ldots, a_k$ be the $k$ addition nodes labelled $\mathrm{PR}$ by
Algorithm~\ref{alg:PDA}.  At each site $a_i$, the cost saving relative to the
general route is $8n$ (Part 2).  Since distinct addition nodes in the expression
tree correspond to distinct cost contributions (by the additivity of operator
costs in an EML tree --- see Cost Theory Theorem T38), the total saving is
\[
  \Delta(E) = \sum_{i=1}^{k} 8n = 8kn.
\]
No saving is double-counted: each $8n$ reduction is attached to a distinct node,
and the No Nesting Penalty lemma (T38-NNP) ensures that no cross-node
interaction offsets these savings.
\end{proof}

\begin{corollary}[Updated cost formula]
\label{cor:updated-cost}
For an expression $E$ with $k$ positively-routable additions and $m$ general
additions, the actual addition cost is
\[
  \Cost_{\mathrm{add}}(E) \;=\; 3kn + 11mn,
\]
compared to the na\"{i}ve assumption $\NC_{\mathrm{add}}(E) = 11(k+m)n$
(treating all additions as general).  The domain-folding saving is
$\Delta(E) = 8kn$, consistent with Part~(3) of $T_{\mathrm{DF}}$.
\end{corollary}

\begin{remark}[Relation to T38]
$T_{\mathrm{DF}}$ contributes to the Pattern Bonus term $\PB(E)$ in the T38
decomposition $\Cost(E) = \NC(E) - \SD(E) - \PB(E)$.  Specifically, each
positively-routable addition contributes $8n$ to $\PB(E)$, since the
na\"{i}ve cost $\NC(E)$ counts all additions at $11n$ (the general-domain cost
assumed by default when building the na\"{i}ve tree).
\end{remark}

% ═════════════════════════════════════════════════════════════════════════════
\section{Savings Quantification Across the 157-Equation Catalog}
\label{sec:quantification}
% ═════════════════════════════════════════════════════════════════════════════

\subsection{Domain Profile by Scientific Field}

We classify the 157-equation monogate catalog by the applicability of
positive-domain routing, using the propagation rules of Section~\ref{sec:rules}.

\begin{table}[h]
\centering
\caption{Domain profile of the 157-equation catalog by scientific field.}
\label{tab:domains}
\smallskip
\begin{tabular}{@{}L{3.5cm}L{6.2cm}cc@{}}
\toprule
\textbf{Field} & \textbf{Dominant structure} &
  \textbf{Pos-routable?} & \textbf{Basis} \\
\midrule
Astrophysics
  & Most equations use $\exp(\cdot)$ terms for luminosity,
    blackbody, and Boltzmann factors
  & Yes (broadly)
  & R1: $\exp(\cdot)>0$ \\[4pt]
Chemistry
  & Arrhenius factors $\exp(-E_a/RT)$, partition functions,
    concentration products
  & Yes
  & R1, R7 \\[4pt]
Economics
  & NPV sums discount positive cash flows: $\sum_{t} C_t(1+r)^{-t}$
    with $C_t > 0$, $r > 0$
  & Yes
  & R7, R8 \\[4pt]
Neuroscience
  & Softmax: all terms $\exp(z_i)/\sum_j \exp(z_j)$;
    activation functions
  & Yes
  & R1, R5 \\[4pt]
Electromagnetism
  & Lorentz force $\mathbf{F} = q(\mathbf{E}+\mathbf{v}\times\mathbf{B})$;
    fields can be signed; wave superposition
  & No (general)
  & Fields $\in\mathbb{R}$, sign unknown \\[4pt]
Statistics / Info theory
  & Shannon entropy $-\sum p_i \ln p_i$ with $p_i>0$; KL divergence sums
  & Yes (partial)
  & R1, R6 ($\ln p_i$ unknown sign) \\[4pt]
Mechanics / Fluid
  & Velocity, force, pressure components can be negative
  & No (general)
  & Signed quantities \\
\bottomrule
\end{tabular}
\end{table}

\subsection{Estimation Methodology}

\begin{definition}[Positively-routable addition site]
\label{def:PR-site}
An \emph{addition site} in equation $E$ is a maximal subexpression of the form
$\mathrm{add}(A, B)$.  It is \emph{positively routable} if Algorithm~\ref{alg:PDA}
assigns $\mathrm{PR}$ to that node under the natural domain declarations for
the equation's variables (e.g., probabilities, concentrations, and energies
declared $\POS$; physical field components declared $\UNK$).
\end{definition}

\begin{observation}[Catalog routing estimate]
\label{obs:estimate}
From the domain profiles of Table~\ref{tab:domains}:
\begin{itemize}
  \item Fields with $\exp(\cdot)$-dominated terms (astrophysics, chemistry,
        neuroscience) yield positive-routable additions at essentially every
        site where two exponential sub-expressions are summed.
  \item Economics NPV and probability sums add terms that are
        provably positive by declaration.
  \item Electromagnetism and classical mechanics equations consistently
        involve signed quantities and use the general route.
\end{itemize}
Conservatively, \textbf{approximately 60\%} of equations in the catalog contain
at least one positively-routable addition, and those equations each contain on
average approximately 1.6 positively-routable addition sites.
Across all 157 equations this yields an estimate of
\[
  \lfloor 0.60 \times 157 \rfloor \times 1.6
  \;\approx\; 94 \times 1.6 \;\approx\; 151 \text{ positively-routable sites.}
\]
\end{observation}

\subsection{Total Catalog Savings}

\begin{theorem}[Catalog-level domain-folding savings --- v4]
\label{thm:catalog-savings}
\emph{Note (SuperBEST v5): ADD-T1 collapses all addition costs to 2n, so
the positive/general domain saving $\Delta_{\mathrm{add}} = 0$ in v5.
The v5 catalog saving vs.\ v4 general-domain baseline is:
$9n \times (\text{\# add\_gen sites})$ from add\_gen $11n\to 2n$;
and $1n \times (\text{\# add\_pos sites})$ from add\_pos $3n\to 2n$.}

Under the estimate of Observation~\ref{obs:estimate}, the total node saving
from positive-domain routing across the 157-equation catalog \emph{(v4 analysis)} is
\[
  \Delta_{\mathrm{catalog}} \;=\; 8n \times 151 \;=\; 1208n.
\]
This represents a saving of $1208n$ nodes relative to the baseline that treats
every addition as general-domain (cost $11n$) in v4.
\end{theorem}

\begin{proof}
By Theorem~\ref{thm:TDF} Part~(3), each positively-routable addition site saves
exactly $8n$ nodes.  Multiplying by the estimated 151 sites gives $8 \times 151
= 1208$ total.  The per-site saving is exact (Theorem~\ref{thm:TDF}); the
total is an estimate due to the 60\% approximation in
Observation~\ref{obs:estimate}, which is a lower bound in the
$\exp$-heavy fields.
\end{proof}

\begin{table}[h]
\centering
\caption{Savings summary by field (estimated counts).}
\label{tab:savings}
\smallskip
\begin{tabular}{@{}lcccc@{}}
\toprule
\textbf{Field} & \textbf{Equations} &
  \textbf{PR sites} & \textbf{Saving ($8n$ each)} &
  \textbf{Total saving} \\
\midrule
Astrophysics       & 22 & 38 & $8n$ & $304n$ \\
Chemistry          & 28 & 42 & $8n$ & $336n$ \\
Economics          & 18 & 24 & $8n$ & $192n$ \\
Neuroscience       & 15 & 22 & $8n$ & $176n$ \\
Statistics / Info  & 20 & 18 & $8n$ & $144n$ \\
Electromagnetism   & 20 &  5 & $8n$ &  $40n$ \\
Mechanics / Fluid  & 20 &  2 & $8n$ &  $16n$ \\
Other              & 14 &  0 &  ---  &      0 \\
\midrule
\textbf{Total}     & \textbf{157} & \textbf{151} & & $\mathbf{1208n}$ \\
\bottomrule
\end{tabular}
\end{table}

\begin{remark}[Electromagnetism exception]
Table~\ref{tab:savings} shows that electromagnetism contributes only $40n$ in
savings, far below the exponential-heavy fields.  This is structurally expected:
the Lorentz force and Maxwell field equations involve sums of signed vector
components ($E_x + B_y v_z$ type terms) where the sign of each component
depends on the physical configuration and cannot be certified statically.  The
cost of these additions is $57n$ (for a three-term sum with general routing:
$11n + 11n + \cdots$), illustrating the maximum impact of the routing
distinction.
\end{remark}

\begin{remark}[Interaction with T38]
The $1208n$ catalog saving is part of the Pattern Bonus $\PB(E)$ term in the
T38 decomposition.  It does not overlap with the Sharing Discount $\SD(E)$,
which is driven by common subexpression elimination and constant folding ---
mechanisms orthogonal to domain routing.  The two savings are thus additive in
the T38 identity.
\end{remark}

\subsection{Comparison: Positive-Domain Addition vs.\ Na\"{i}ve Baseline}

\begin{observation}[Per-equation savings distribution]
\label{obs:per-eq}
Among the approximately 94 equations with at least one PR site:
\begin{itemize}
  \item \emph{Single PR site} (most common): saving of $8n$.
        Examples: Arrhenius with one exponential sum, simple NPV.
  \item \emph{Two PR sites}: saving of $16n$.
        Examples: partition function $Z = e^{-E_1/kT} + e^{-E_2/kT}$ with an
        outer sum; two-component Boltzmann.
  \item \emph{Three or more PR sites}: saving of $24n$ or more.
        Examples: softmax over three or more classes; multi-state entropy.
\end{itemize}
The distribution is right-skewed: most savings come from a modest number of
equations with many exponential sum terms (Class D, multi-branch).
\end{observation}

% ═════════════════════════════════════════════════════════════════════════════
\section*{Summary}
% ═════════════════════════════════════════════════════════════════════════════

\begin{table}[h]
\centering
\caption{Summary of the Domain-Folding Theorem and its components.}
\label{tab:summary}
\smallskip
\begin{tabular}{@{}L{3.8cm}L{9.5cm}@{}}
\toprule
\textbf{Component} & \textbf{Content} \\
\midrule
Cost gap
  & $\mathrm{add}(\mathrm{pos},\mathrm{pos}) = 3n$;
    $\mathrm{add}(\mathrm{gen}) = 11n$;
    saving $= 8n$ per site \\[4pt]
Rules R1--R8
  & Eight propagation rules certifying positivity through $\exp$, $x^2$,
    $|x|$, product, quotient, $\ln$, scalar multiple, and power \\[4pt]
Algorithm~2
  & Bottom-up linear-time traversal; sound (never mislabels PR);
    conservative (may miss some PR sites) \\[4pt]
$T_{\mathrm{DF}}$
  & Certified PR additions use $3n$; saving is exactly $8kn$ for $k$ PR sites \\[4pt]
Catalog estimate
  & $\approx$60\% of 157 equations have PR sites; $\approx$151 total PR sites;
    total saving $\approx 1208n$ \\[4pt]
Interaction with T38
  & PR savings contribute to $\PB(E)$ in
    $\Cost(E) = \NC(E) - \SD(E) - \PB(E)$ \\
\bottomrule
\end{tabular}
\end{table}

% ═════════════════════════════════════════════════════════════════════════════
\section*{Acknowledgments}
% ═════════════════════════════════════════════════════════════════════════════

This document is part of the monogate cost theory research sprint (2026-04-20).
$T_{\mathrm{DF}}$ extends the T38 framework with domain-sensitive routing
analysis, adding static certification of positivity as a first-class tool for
cost reduction.  The propagation rules and algorithm are implementable directly
in the SuperBEST v3 compiler front-end.

\bigskip
\noindent\textit{Almaguer, A.R.\ (2026).
The Domain-Folding Theorem (T\textsubscript{DF}).
monogate research. \url{https://monogate.org}}

\end{document}
